%% ============================================================================
%% Fichier  : annexes/Ann_B_outils.tex
%% Titre    : Référentiel des Outils
%% Version  : 1.0 -- novembre 2026 (INF3226, d'après ACSI)
%% ============================================================================
%% Outils effectivement mobilisés dans les chapitres et les labs ; les noms
%% de logiciels commerciaux sont des exemples, à vérifier au moment du choix
%% (Chapitre 19).
%% ============================================================================

\chapter{Référentiel des Outils}
\label{ann:outils}

Cette annexe rassemble les outils mobilisés dans le cours et dans les labs
de NDAP SUKULU. Les noms de logiciels commerciaux sont des exemples~: les
offres évoluent, et le choix se fait selon les six critères du
Chapitre~\ref{chap:audit-continu}.

\section*{Intégrité et preuve}

\begin{tableaucours}{L{3.4cm}L{5.2cm}L{3.6cm}}{%
    \tableauHeader{Outil} &
    \tableauHeader{Usage} &
    \tableauHeader{Chapitre, lab}
}
\texttt{sha256sum} & Calculer et vérifier l'empreinte d'une extraction
(\texttt{sha256sum -c MANIFESTE\_SHA256.txt}) & Ch.~\ref{chap:preuve}~;
tous les labs \\
Lettre d'extraction & Source, requête, date, volumes, empreinte &
Ch.~\ref{chap:preuve}, ASI-LAB-07 \\
Registre de conservation & Chaque accès et chaque remise d'une pièce &
Ch.~\ref{chap:audit-preuve} \\
\end{tableaucours}

\section*{Analyse de données}

\begin{tableaucours}{L{3.4cm}L{5.2cm}L{3.6cm}}{%
    \tableauHeader{Outil} &
    \tableauHeader{Usage} &
    \tableauHeader{Chapitre, lab}
}
Tableur (LibreOffice Calc) & Filtres, tableaux croisés, \texttt{NB.SI.ENS},
\texttt{RECHERCHEV}, normalisation des références & ASI-LAB-01 à ASI-LAB-10 \\
SQL & Doublons, équilibre, seuils, séquences, rapprochements
(\texttt{GROUP BY}, \texttt{HAVING}, \texttt{LEFT JOIN}) &
Ch.~\ref{chap:audit-applications}, ASI-LAB-07 \\
Langage de script (Python et ses bibliothèques d'analyse) & Tests
reproductibles sur de grands volumes, contrôle mensuel, loi de Benford,
simulation FAIR & Ch.~\ref{chap:gestion-risques}, \ref{chap:fraude},
\ref{chap:audit-continu} \\
Logiciels d'analyse pour auditeurs (par exemple ACL, IDEA) & Tests
intégrés, journal des traitements, sans programmation &
Ch.~\ref{chap:audit-continu} \\
\end{tableaucours}

\section*{Gestion de la mission et du risque}

\begin{tableaucours}{L{3.4cm}L{5.2cm}L{3.6cm}}{%
    \tableauHeader{Outil} &
    \tableauHeader{Usage} &
    \tableauHeader{Chapitre, lab}
}
Classeur O4 & Décomposition, fiches d'assertion, cotation, $\mu$, $U$, $H$
& Ch.~\ref{chap:decomposition} à~\ref{chap:conclure}, ASI-LAB-04, ASI-LAB-05 \\
Registre des risques & Probabilité, impact, score, appétence &
Ch.~\ref{chap:risque}, ASI-LAB-01 \\
Méthodes d'appréciation (EBIOS RM, MEHARI, OCTAVE, FAIR) & Scénarios,
ateliers, quantification & Ch.~\ref{chap:gestion-risques} \\
Déclaration d'applicabilité & Mesures retenues et exclues, justification
& Ch.~\ref{chap:audit-smsi} \\
Matrice des droits & Paires incompatibles, comptes sans titulaire &
Ch.~\ref{chap:controle-interne}, ASI-LAB-02 \\
Plateformes de gouvernance, risques et conformité & Cartographie,
auto-évaluations, suivi des recommandations & Ch.~\ref{chap:audit-continu} \\
\end{tableaucours}

\section*{Tests techniques}

\begin{tableaucours}{L{3.4cm}L{5.2cm}L{3.6cm}}{%
    \tableauHeader{Outil} &
    \tableauHeader{Usage} &
    \tableauHeader{Chapitre, lab}
}
Exports de configuration & Règles des pare-feu, services, comptes,
correctifs, journalisation & Ch.~\ref{chap:audit-technique}, ASI-LAB-08 \\
Scanners de ports et de vulnérabilités & Uniquement après note
d'autorisation signée, dans le périmètre et la fenêtre convenus &
Ch.~\ref{chap:audit-technique}, \ref{chap:intrusion} \\
Référentiels de test (OWASP ASVS, Top~10, NIST SP~800-115, CVSS) &
Construire le plan de test, noter la gravité & Ch.~\ref{chap:intrusion} \\
Repères CIS & Configurations de référence des systèmes &
Ch.~\ref{chap:audit-technique} \\
\end{tableaucours}

\begin{boxalert}
\textbf{Rappel.} Aucun scan n'est réalisé dans les labs du cours~: le
ASI-LAB-08 fournit des exports figés. Un test technique sur un système réel
exige une note d'autorisation écrite (loi n°~2010/012, art.~65, 69, 70 et~86~;
Chapitre~\ref{chap:intrusion}).
\end{boxalert}

\section*{Le bureau de NDAP SUKULU}

Les labs s'exécutent dans un bureau d'analyste virtuel. Les pièces sont
remises dans \texttt{\textasciitilde/preuves}, en lecture seule, avec leur
manifeste d'empreintes~; une copie de travail se trouve dans
\texttt{\textasciitilde/travail}. Le bureau propose un terminal et une suite
bureautique. Toute l'activité dans le terminal est journalisée.
